\section*{الفصل \ref{ch:g8:electric-current-effects} --- التيار الكهربائي وآثاره}

\begin{solution}{exo:g8:electric-current-effects:1}
الحراري (المحمصة، والغلّاية، وخيط \omterm{def:g3:first-electric-circuit:bulb}{المصباح})؛ والمغناطيسي
(\omterm{def:g8:electric-current-effects:electromagnet}{المغناطيس الكهربائي}، \omterm{ex:g6:circuits-and-safety:motor}{والمحرّك}، وجرس الباب)؛ والكيميائي (الطلي
الكهربائي، وشطر الماء، وشحن البطاريات).
\end{solution}

\begin{solution}{exo:g8:electric-current-effects:2}
حراري؛ فمغناطيسي؛ فكيميائي؛ فحراري؛ فمغناطيسي.
\end{solution}

\begin{solution}{exo:g8:electric-current-effects:3}
أن التيار يخلق مغناطيسية حول سلكه --- وهو أول جسر بين الكهرباء
والمغناطيسية. وعكس التيار يعكس الرجفة.
\end{solution}

\begin{solution}{exo:g8:electric-current-effects:4}
يمكن إطفاؤه وتشغيله وقلب قطبيه متى شئت --- ويمكن تجريع قوّته
(لفّات أكثر، وتيار أقوى). \omterm{def:g1:magnets:magnet}{والمغناطيس} الحجري لا يفعل شيئًا من
ذلك.
\end{solution}

\begin{solution}{exo:g8:electric-current-effects:5}
اللفّات العارية المتلامسة ستتيح للتيار أن يختصر عبر اللفّ بدل
أن يجول في الوشيعة --- فالعزل يُلزمه الطريق المغنطِط الطويل.
والوشيعة وحدها لا تكاد تقاوم المسير: فإذا تُركت موصولة كانت
\omterm{def:g7:short-circuits-safety:device}{دارة قصيرة} تقريبًا، تسخّن الوشيعة \omterm{def:g3:first-electric-circuit:battery}{والبطارية} --- \omterm{prop:g8:electric-current-effects:three}{والأثر الحراري}
مرافق.
\end{solution}

\begin{solution}{exo:g8:electric-current-effects:6}
بفتح قاطعة: فالتيار زائل، فالمغناطيسية زائلة، فالسيارة ساقطة.
أما \omterm{def:g1:magnets:magnet}{مغناطيس} حجري بتلك \omterm{def:g2:pushes-and-pulls:force}{القوة} فسيمسك إلى الأبد --- ولاحتاج إطلاق
السيارة إلى آلات تنتزعها انتزاعًا.
\end{solution}

\begin{solution}{exo:g8:electric-current-effects:7}
قيد الاستعمال: \omterm{def:g5:everyday-energy:energy}{الطاقة} الكيميائية المخزونة تصير \omterm{def:g5:everyday-energy:energy}{طاقة} كهربائية
تسوق المسير (فتنتهي ضوءًا وحركة ودفئًا). وعلى الشاحن: \omterm{def:g5:everyday-energy:energy}{طاقة}
كهربائية مساقة بالمقلوب عبر الخلية تعيد بناء المخزون
الكيميائي.
\end{solution}

\begin{solution}{exo:g8:electric-current-effects:8}
لأن المسير نفسه غير مرئي: فلا يفضحه إلا أفعاله. فالتدفئة
والرجفة والفقاعات هي توقيعات التيار الوحيدة --- وآلة الفصل
التالي تصقل الرجفة المغناطيسية في قراءة مرقّمة.
\end{solution}

\begin{solution}{exo:g8:electric-current-effects:9}
يُضغط الزرّ: فيسري التيار، فيخطف \omterm{def:g8:electric-current-effects:electromagnet}{المغناطيس الكهربائي} المطرقة،
فرنّة --- لكن أرجحة المطرقة تفتح تلامسًا في تغذيتها هي؛
فتموت المغناطيسية، فيعيدها النابض، فيعاد صنع التلامس؛ فيعود
التيار، فخطفة أخرى --- وهكذا مرات كثيرة في الثانية. والإمساك
مستحيل بحكم التصميم: فالخطفة تقطع سببها دائمًا --- والخشخشة
\emph{هي} الرنين.
\end{solution}

\begin{solution}{exo:g8:electric-current-effects:10}
تلامسات منظّم \omterm{def:g5:heat-and-insulation:heat}{الحرارة} الواهية لا تحمل إلا تيار الوشيعة الصغير؛
ومغناطيسية الوشيعة تغلق تلامسات متينة تحمل تيار المسخّن النهم
في حلقة منفصلة. فالصغير يأمر الكبير، والدارتان لا تتلامسان
أبدًا --- فاليد الرقيقة تجذب \omterm{def:g3:levers-and-scales:lever}{رافعة}، لا الحمل.
\end{solution}

\begin{solution}{exo:g8:electric-current-effects:11}
البروتوكول: أمسك \omterm{def:g4:magnets-and-compass:compass}{البوصلة} مسطّحة على الجدار، بعيدًا عن الأثاث
الحديدي؛ وامسح ببطء؛ وحيث ترتجف الإبرة عند تشغيل الجهاز
المشتبه به وإطفائه، يمرّ سلك حامل للتيار. والرجفة القوية تعني
تيارًا قريبًا أو قويًّا. والحدّ الأمين: الإبر تجيب أيضًا
الأنابيبَ الحديدية والمغانطَ والأرضَ نفسها --- فلا تثبت تيارًا
إلا رجفة تتبع تشغيل \omterm{ex:g3:first-electric-circuit:switch}{القاطعة} وإطفاءها.
\end{solution}

\begin{solution}{exo:g8:electric-current-effects:12}
التيار يجعل الوشيعة مغناطيسًا؛ فتتجاذب أقطاب الوشيعة والإطار
المختلفة وتتنافر المتشابهة، فتدور الوشيعة نحو التطابق. وبدون
العكس كانت ستبلغ التطابق وتقف --- نصف دورة، ثم ارتجافة إلى
السكون: مغلاق باب لا محرّك. أما عكس التيار عند لحظة التطابق
فيبدّل قطبَي الوشيعة: فجذب الأمس صار تنافرًا، وصار التطابق أسوأ
موضع، ووجب على الوشيعة أن تواصل مطاردة هدف ينقلب هاربًا كل نصف
دورة --- مطاردة أبدية، أي: دوران.
\end{solution}

\begin{solution}{pb:g8:electric-current-effects:1}
\textbf{1.} العروض الثلاثة كلها على \emph{التسلسل} في حلقة
واحدة مع \omterm{ex:g3:first-electric-circuit:switch}{القاطعة}: فمسير واحد يخدم عندئذٍ السلك الدافئ وسلك
\omterm{def:g4:magnets-and-compass:compass}{البوصلة} والجرّة بالتناوب --- قاطعة واحدة، وثلاثة آثار.
\textbf{2.} قانون التسلسل: التيار \emph{نفسه} يعبر السلك
الغليظ والرفيع سواءً. والمختلف هو الطريق لا غير --- فالسلك
الرفيع يعاني المسير أكثر فيدفأ: الجرعة نفسها، والمريض مختلف.
\textbf{3.} اعكس \omterm{def:g3:first-electric-circuit:battery}{البطارية} \emph{ومرّر} السلك من جهة \omterm{def:g4:magnets-and-compass:compass}{البوصلة}
الأخرى (أو اقلب \omterm{def:g4:magnets-and-compass:compass}{البوصلة}): فعكسان يلغي أحدهما الآخر، وتبقى
الرجفة شرقًا.
\textbf{4.} أن مكوّنَي الماء حاضران بنسبة اثنين إلى واحد ---
\omterm{def:g3:solids-liquids-gases:gas}{فالغاز} من إحدى الرصاصتين ضِعف \omterm{def:g3:solids-liquids-gases:gas}{الغاز} من الأخرى.
\textbf{5.} أزِل التلامس القاطع لنفسه (القطعة التي تفتحها
أرجحة المطرقة): فيمسك \omterm{def:g8:electric-current-effects:electromagnet}{المغناطيس الكهربائي} عندئذٍ مطرقة اللبّاد
على الجرس في صمت --- ولا يسمع الزائر إلا ``دُمّ'' واحدة
ناعمة، ولا شيء بعدها.
\textbf{6.} عند الصندوق يفتح المشغّل قاطعة الوشيعة: فتزول
المغناطيسية، ويسقط الحديد. أما الألمنيوم فيواصل الركوب لأن
المغانط --- حجرية كانت أم من صنع التيار --- لا تمسك إلا عائلة
الحديد، والألمنيوم لم يجبها قطّ.
\textbf{7.} أن أقطاب \omterm{def:g1:magnets:magnet}{المغناطيس} الذي يصنعه التيار تتبع اتجاه
التيار: فاعكس أحدهما ينعكس الآخر.
\textbf{8.} \omterm{def:g3:first-electric-circuit:battery}{البطارية} تسوق الطلي \emph{بإنفاق} مخزونها
الكيميائي؛ وإعادة بناء مخزون تقتضي \omterm{def:g5:everyday-energy:energy}{طاقة} مساقة إليه من الخارج.
\omterm{def:g3:first-electric-circuit:battery}{والبطارية} التي تعيد ملء نفسها من العمل نفسه الذي تدفع ثمنه هي
آلة الأبد ثانية --- ففاتورة \omterm{def:g5:everyday-energy:energy}{الطاقة} لا تستطيع أن تدفع نفسها.
\textbf{9.} الإمساك: تُغلق \omterm{ex:g3:first-electric-circuit:switch}{القاطعة}، فتتمغنط الوشيعة، فتُمسَك
الحلية. والأرجحة: التيار محفوظ، فالقبضة محفوظة --- \omterm{prop:g8:electric-current-effects:three}{والأثر
المغناطيسي} في الخدمة. والإطلاق: تُفتح \omterm{ex:g3:first-electric-circuit:switch}{القاطعة} فوق المزلق،
فتموت المغناطيسية، فتسقط الجائزة.
\textbf{10.} أن \omterm{def:g8:electric-current-effects:electromagnet}{المغناطيس الكهربائي} يحيا بتياره: \omterm{def:g3:first-electric-circuit:battery}{فبطارية}
هزيلة، فمسير أضعف، فقبضة أضعف --- وثمن القابلية للإطفاء هو
التبعية. (\omterm{def:g1:magnets:magnet}{والمغناطيس} الحجري لا يهزل أبدًا --- ولا يترك
أبدًا.)
\textbf{11.} لفّات أكثر، \omterm{def:g1:magnets:magnet}{فمغناطيس} أقوى عند التيار نفسه: فقبضة
أمتن. والجرعة الأخرى هي \omterm{def:g2:pushes-and-pulls:force}{قوة} التيار نفسه --- \emph{الشدّة}،
المقيسة في الفصل التالي.
\textbf{12.} مثلًا: ``الحراري: \omterm{def:g7:short-circuits-safety:fuse}{المصهر}، يموت أولًا وبأمانة.
والمغناطيسي: لعبة \omterm{def:g3:levers-and-scales:lever}{الرافعة}، تمسك وتطلق عند الأمر. والكيميائي:
جرّة الطلي، تفضّض الملاعق بصبر \omterm{def:g3:first-electric-circuit:battery}{بطارية}.''
\end{solution}
